%
% Report finale - Laboratorio di Ingegneria Informatica
% Layout a colonna singola (adattato dal template ACL 2018 richiesto dalla consegna)
%
\documentclass[11pt]{article}
\usepackage[a4paper, margin=2.5cm]{geometry}
\usepackage{times}
\usepackage[T1]{fontenc}
\usepackage[utf8]{inputenc}
\usepackage{url}
\usepackage{titlesec}
\usepackage{parskip}
\usepackage{microtype}

\sloppy
\titleformat{\section}{\large\bfseries}{\thesection}{0.6em}{}
\titlespacing*{\section}{0pt}{1.4em}{0.6em}

\title{\textbf{Minerva: una Pipeline di Parsing e Valutazione Automatica \\ per Sorgenti Web Eterogenee}}
\author{
  Laura De Michelis \quad -- \quad Matricola 2119887 \quad -- \quad \texttt{\small demichelis.2119887@studenti.uniroma1.it} \\[2pt]
  Andrea Giuliotti \quad -- \quad Matricola 2107814 \quad -- \quad \texttt{\small giuliotti.2107814@studenti.uniroma1.it} \\[2pt]
  Alessandro Di Nitto \quad -- \quad Matricola 2109155 \quad -- \quad \texttt{\small dinitto.2109155@studenti.uniroma1.it} \\[6pt]
  Sapienza Università di Roma
}
\date{}

\begin{document}
\maketitle

\section{Obiettivo del Progetto}

Con questo progetto abbiamo costruito una pipeline completa che scarica
pagine da quattro domini web diversi (Wikipedia in italiano, AppleVis,
Basketball-Reference e OndaRock), ne estrae il testo utile, lo confronta con
un Gold Standard scritto da noi e lo fa valutare anche da un LLM locale. Tutto
il sistema è esposto tramite API REST e una Web UI, e gira in quattro
container Docker (backend, frontend, database, LLM).

\section{Organizzazione del Codice}

Il backend (FastAPI, porta 8003) sta in \texttt{backend/src/} ed è diviso in
quattro parti. In \texttt{parsers/} c'è una classe astratta
\texttt{BaseDomainParser} che si occupa di tutto quello che è comune a ogni
sito (aprire il browser di Crawl4AI, scaricare l'HTML, riparsarlo da una
stringa già salvata), mentre ogni dominio ha la sua sottoclasse con il
selettore CSS giusto, le esclusioni per togliere il rumore e la pulizia
finale del Markdown -- abbiamo capito quali selettori usare guardando a
mano l'HTML di ogni sito. Un \texttt{ParserFactory} sceglie il parser giusto
in base al dominio dell'URL. In \texttt{evaluation/} c'è tutto quello che
serve a ripulire il Markdown prima del confronto, calcolare le metriche e
parlare con Ollama per il giudizio dell'LLM. In \texttt{db/} c'è la
connessione a MariaDB, lo schema e le funzioni per leggere/scrivere, incluso
lo script che popola il database all'avvio. In \texttt{models/} ci sono gli
schemi Pydantic di input/output di ogni endpoint. Il frontend (FastAPI +
Jinja2, porta 8004) ha quattro pagine (home, parser \& evaluation, gold
standard builder, stats) e parla \emph{solo} con le API REST del backend,
mai direttamente col database.

\section{Valutazione Parser}

Come richiesto abbiamo usato \texttt{token\_level\_eval} (precision, recall,
F1 sui token separati da spazio, dopo aver tolto il Markdown da entrambi i
testi), e in più abbiamo aggiunto ROUGE-1: a differenza della prima metrica,
che guarda solo se una parola c'è o no, ROUGE-1 conta anche quante volte
compare, quindi si accorge anche se il parser ripete o salta pezzi di testo.

Gli F1 medi che abbiamo ottenuto sulle 10 pagine di ogni dominio sono:
\texttt{it.wikipedia.org} 0.90, \texttt{www.applevis.com} 0.97,
\texttt{www.basketball-reference.com} 0.90, \texttt{www.ondarock.it} 0.99.
Per basketball-reference.com abbiamo deciso di far coprire al Gold Standard
solo il blocco anagrafico in alto alla pagina (\texttt{\#meta}: ruolo,
squadra, dati di nascita, statistiche riassuntive), e non le tabelle enormi
di statistiche partita per partita che seguono, perché quelle sono dati
veri ma non ``testo informativo di sintesi'' nel senso che chiede la
consegna.

Durante la costruzione e il controllo del Gold Standard abbiamo trovato due
problemi interessanti. Il primo: il parser HTML della libreria standard di
Python si fermava senza dire niente a metà di una pagina con un markup un
po' irregolare, tagliando via il 95\% del contenuto reale -- ce ne siamo
accorti confrontando la lunghezza che ci aspettavamo con quella che
uscivano, e abbiamo risolto passando a un parser HTML più permissivo. Il
secondo: anche quando l'estrazione è sostanzialmente giusta, piccoli residui
di formattazione Markdown (uno spazio in più o in meno prima dei due punti,
parole attaccate ai bordi di un grassetto o un corsivo dopo aver tolto i
simboli) fanno perdere qualche punto di F1: è un limite di come si
tokenizza il testo (a spazi), più che un vero difetto del parser -- l'abbiamo
verificato controllando a mano i token che non combaciavano tra estrazione e
Gold Standard.

\section{Valutazione dei Judge}

Come LLM Judge abbiamo scelto \texttt{llama3.2:3b} (Meta, circa 2GB di RAM)
tra i cinque modelli proposti, principalmente perché è il più veloce su CPU
e la macchina su cui abbiamo sviluppato non è potentissima: una singola
valutazione ci mette da circa 1 a circa 4 minuti, a seconda che il modello
sia già caricato in RAM o debba essere ricaricato da disco. Il prompt chiede
una risposta in JSON con \texttt{score} (da 1 a 5) e \texttt{feedback}, con
temperatura bassa (0.1) per avere un giudizio il più possibile coerente a
parità di input. Per il parsing della risposta abbiamo previsto quattro
livelli di fallback in cascata (JSON diretto, poi il primo blocco
\texttt{\{...\}} che si trova nel testo, poi i singoli campi via regex,
infine una cifra isolata), e solo se proprio niente funziona ricadiamo su un
punteggio neutro.

Una cosa che ci ha colpito facendo i test: su tutte e 40 le pagine del Gold
Standard il punteggio è quasi sempre 4/5, indipendentemente dal vero F1
(che va da 0.85 a 0.99). Per capire se fosse un bug abbiamo fatto un test di
controllo dandogli un testo completamente scollegato dal Gold Standard, e
lì il punteggio è sceso correttamente a 2/5, con un feedback che citava
proprio l'argomento sbagliato -- quindi il modello legge davvero il
contenuto, non restituisce sempre lo stesso numero. Semplicemente sembra che
riservi il 5/5 quasi solo a corrispondenze testuali identiche, e tratti
qualunque piccola imperfezione (anche minima) allo stesso modo. È un
comportamento abbastanza noto dei modelli piccoli usati come giudici
automatici a bassa temperatura, non un difetto della nostra pipeline.

Per curiosità abbiamo provato anche altri modelli oltre a
\texttt{llama3.2:3b}: \texttt{qwen3:4b} (non compatibile con il nostro
codice così com'è, perché restituisce il ragionamento in un campo
\texttt{thinking} separato invece che in \texttt{response}),
\texttt{phi4-mini} (risultati peggiori, punteggi ancora meno informativi),
\texttt{gemma3:e2b} (7GB, troppo pesante: il caricamento del modello
saturava la RAM e faceva andare in timeout il browser usato per lo
scraping) e \texttt{ministral-3:3b} (3GB, dimensione simile a
\texttt{llama3.2:3b}, ma con generazione a circa 1 token al secondo sulla
nostra CPU: una singola valutazione ha impiegato oltre 7 minuti e la
richiesta successiva ha fatto crashare il processo interno di Ollama).
Questi test confermano che il collo di bottiglia è l'hardware della
macchina di sviluppo (CPU senza GPU dedicata, RAM condivisa tra Docker,
Ollama e il browser per lo scraping) più che la scelta del modello, e che
\texttt{llama3.2:3b} resta il compromesso più stabile tra i cinque
proposti.

Abbiamo anche provato a migliorare il prompt del judge dando una rubrica
esplicita (criteri diversi per ogni punteggio da 1 a 5, invece della
richiesta generica) su tre pagine reali con F1 molto diverso tra loro
(1.00, 0.92, 0.81): il punteggio restituito è stato 4/5 in tutti e tre i
casi, sia col prompt vecchio che con quello nuovo. Il feedback testuale
cambiava correttamente da un caso all'altro (segno che il modello legge
davvero i testi), ma il punteggio numerico no. Consideriamo quindi questo
comportamento un limite di calibrazione del modello più che un problema
risolvibile lato nostro senza introdurre un aggiustamento manuale del
punteggio, cosa che abbiamo deciso di non fare per non falsare il giudizio
dell'LLM.

\section{Organizzazione del DB}

Nel database MariaDB abbiamo le due tabelle obbligatorie
\texttt{web\_resources} e \texttt{gold\_standard}, collegate uno-a-uno con
una chiave esterna su \texttt{url} (con cascade in eliminazione), più due
tabelle nostre: \texttt{evaluations} (i risultati di
\texttt{token\_level\_eval} e ROUGE-1 per ogni pagina) e
\texttt{llm\_judgments} (punteggio, feedback e nome del modello usato per il
giudizio). Tutte le query passano da un livello di repository e sono
parametrizzate. All'avvio del backend, uno script di seed legge tutti i file
\texttt{gs\_data/*.json} e riempie \texttt{web\_resources} e
\texttt{gold\_standard} con un upsert, così l'HTML grezzo di ogni pagina
resta sempre disponibile per rifare il parsing e confrontarlo col
\texttt{gold\_text} quando vogliamo, senza dover riscaricare niente.



\end{document}
